\section{Choosing and validating a trusted region from unlabeled data}
\label{sec:unlabeled-trusted-region}

A trusted region need not be specified before seeing any data. Independent
unlabeled sample splitting supplies a simple rigorous version of this idea.
This is a standard statistical construction, not a new sample-splitting
principle. Its role here is to make the coverage assumption in the preceding
robust-prior results independently checkable.

Let $D$ be the test input distribution on $\{-1,1\}^d$. Draw three mutually
independent samples: an unlabeled pilot $P$, an unlabeled validation sample
$V=(X'_1,\ldots,X'_m)$, and the labeled training sample
$S=((X_i,f(X_i)))_{i=1}^n$. Using only $P$, choose any measurable region
$G=G(P)$. The selection rule can be arbitrarily complicated. For example,
$P$ may determine a density level set, a union of neighborhoods of observed
inputs, or a support estimate in a fixed representation. The theorem does
not assert that such regions admit an easy all-input stability certificate.
Define
\[
 \widehat\rho_G=\frac1m\sum_{j=1}^m\mathbf1\{X'_j\notin G\},\qquad
 \overline\rho_G=\min\!\left\{1,\widehat\rho_G+
            \sqrt{\frac{\log(1/\delta_u)}{2m}}\right\}.
\]

\begin{theorem}[An independently learned trusted region]
\label{thm:unlabeled-trusted-region}
Fix $0<\delta_u,\delta_l<1$ and $0<\varepsilon<1$.
Suppose the target has a width-$C$ layered Boolean circuit. For the selected
$G(P)$, form any one of the following priors by conditioning the raw Gaussian
on the indicated condition at every $x\in G(P)$:
\begin{enumerate}
\item The calibrated sign stability condition of
Theorem~\ref{thm:trusted-region}, followed by ordinary training sign
consistency. Let $R(Q)$ denote Gibbs sign error.
\item The numerical stability condition of
Theorem~\ref{thm:numeric-prior}, with signed training outputs greater
than $1/2$. Let $R(Q)$ denote
$\Pr_{w\sim Q,X\sim D}\{f(X)z_w(X)\le1/4\}$.
\item The clipped-score stability condition of
Corollary~\ref{cor:clipped-stability}, with the same training-margin
likelihood and the same test loss as in item 2.
\end{enumerate}
No stability condition is imposed outside $G(P)$, and no additional
all-input clean-margin restriction is imposed. For the clipped condition, fix $0<c\le1/100$ once and for all;
constants in the following bound may depend on $c$. The numerical
condition retains its stated fixed calibration. There is a
polynomial sample bound
\[
 n\ge\operatorname{poly}
       (d,L,C,\varepsilon^{-1},\log(1/\delta_l))
\]
such that, with probability at least $1-\delta_u-\delta_l$ over $P,V,S$,
\[
 R(Q)\le\overline\rho_G+\varepsilon,
\]
simultaneously for all widths $N\ge C$ and scales $0<t\le1$.
The complexity of the pilot selection rule does not appear in this bound.
The Gaussian is conditioned exactly on the selected support condition;
no robust training likelihood or soft reweighting is introduced.
\end{theorem}

\begin{proof}
Condition on the pilot. The set $G$ is now fixed and independent of both
$V$ and $S$. Hoeffding's inequality implies
$D(G^c)\le\overline\rho_G$ with probability at least $1-\delta_u$ over
$V$. The corresponding fixed-region theorem gives
$R(Q)\le D(G^c)+\varepsilon$ with probability at least $1-\delta_l$ over
$S$. Its constants and common finite compression class are uniform over
$G$: in particular, the globally stable target cylinder belongs to the
conditioned support for every possible $G$. A union bound, followed by
averaging over the pilot, proves the result. The sample event supplied by
the fixed-region theorem is simultaneous in width and Gaussian scale.
\end{proof}

For any of these losses, a pointwise posterior correct-answer probability
of at most $2/3$ entails pointwise Gibbs error at least $1/3$.
Consequently the measure of such inputs is at most
$3(\overline\rho_G+\varepsilon)$. For items 2 and 3, ``correct answer''
can mean the stronger event $f(x)z_w(x)>1/4$.
An exact one-sided binomial upper confidence limit can replace Hoeffding.
For example, when the validation sample has no misses, the upper limit is
$1-\delta_u^{1/m}\le\log(1/\delta_u)/m$.

\paragraph{Validation candidates.}
The pilot may produce a catalog of $M$ regions, denoted $G_j(P)$.
Using
\[
 \overline\rho_j=\min\!\left\{1,\widehat\rho_j+
              \sqrt{\frac{\log(M/\delta_u)}{2m}}\right\}
\]
certifies all $M$ coverages simultaneously. One may then choose any single
index $J=J(P,V)$ before seeing $S$, for example the index with the best
coverage among regions admitting a convenient certificate. The theorem
holds with $G_J$ and $\overline\rho_J$. There is no extra $\log M$ cost in
the labeled sample bound: conditional on $P,V$, just one region has been
chosen independently of $S$. The $\log M$ cost is paid in the unlabeled
validation confidence bound. If the region index is instead chosen using
$S$, the additional complexity control of the finite-catalog theorem is
needed. Reusing validation to generate an unbounded number of new regions
also requires further control; the single-region guarantee cannot simply
be applied after that adaptation.

\paragraph{What is still required.}
Coverage validation checks only how often an input lies in $G$. The prior
must still impose stability throughout $G$, simultaneously for every
admissible disturbance. Certifying stability merely on the finite pilot or
validation inputs is not sufficient. In particular, high average certified
robustness of one trained network does not by itself prove a theorem about
the Bayesian posterior. A network-dependent exceptional region can adapt
to the training labels, which is precisely the difficulty exhibited by the
sparse-patch example. The present construction avoids that difficulty by
freezing a shared region before the labeled sample. It modifies the support
of the Gaussian prior; it does not establish the same conclusion for the
unrestricted Gaussian prior.
